\documentclass[12pt,a4paper]{article}
\usepackage{fontspec}\setmainfont{Latin Modern Roman}\setsansfont{DejaVu Sans}
\usepackage{amsmath,amssymb,graphicx,geometry,enumitem,needspace,longtable,booktabs,fancyhdr,hyperref}
\geometry{left=28mm,right=23mm,top=24mm,bottom=24mm}
\hypersetup{unicode=true,colorlinks=true,linkcolor=black,urlcolor=blue,pdftitle={Giải tích §4.9 — Nguyên hàm}}
\setlength{\parindent}{1.2em}\setlength{\parskip}{4pt}\setlength{\headheight}{16pt}
\pagestyle{fancy}\fancyhf{}\fancyhead[L]{\small Chương 4. Ứng dụng của đạo hàm}\fancyhead[R]{\small §4.9}\fancyfoot[C]{\thepage}
\newcommand{\fig}[3][.68]{\begin{center}\begin{minipage}{\linewidth}\centering\includegraphics[width=#1\linewidth]{assets/#2.png}\\{\small #3}\end{minipage}\end{center}}
\newcommand{\ex}[1]{\par\medskip\noindent\textbf{Ví dụ #1.}\ }
\newcommand{\sol}{\par\nopagebreak[4]\noindent\textit{Lời giải.}\ }
\newcommand{\bt}[1]{\par\smallskip\Needspace{5\baselineskip}\noindent\textbf{#1.}\ }
\newcommand{\group}[1]{\par\medskip\noindent\textbf{#1}\par\nopagebreak[4]}
\newcommand{\gint}[1]{\left[\!\left[#1\right]\!\right]}
\allowdisplaybreaks[2]
\DeclareMathOperator{\sech}{sech}
\DeclareMathOperator{\csch}{csch}
\DeclareMathOperator{\arccot}{arccot}
\begin{document}
\section*{4.9 — Nguyên hàm}
\noindent\textit{Bản dịch tiếng Việt — Stewart, Clegg và Watson, Calculus: Early Transcendentals, ấn bản thứ 9. Phạm vi: §4.9, tr. 356–364.}

Nhà vật lý biết vận tốc của một chất điểm có thể muốn tìm vị trí của nó tại một thời điểm. Kỹ sư đo được tốc độ rò rỉ biến thiên của nước từ bể muốn biết lượng nước đã rò trong một khoảng thời gian. Nhà sinh học biết tốc độ tăng của quần thể vi khuẩn có thể muốn suy ra kích thước quần thể trong tương lai. Trong mỗi trường hợp, bài toán là tìm một hàm có đạo hàm là một hàm đã biết.

\subsection*{Nguyên hàm của một hàm số}
Nếu hàm $F$ có đạo hàm là $f$, thì $F$ được gọi là một nguyên hàm của $f$.
\group{Định nghĩa}
Hàm $F$ được gọi là một \textbf{nguyên hàm} của $f$ trên khoảng $I$ nếu $F'(x)=f(x)$ với mọi $x\in I$.

Chẳng hạn, xét $f(x)=x^2$. Nếu nhớ quy tắc lũy thừa, không khó tìm một nguyên hàm: $F(x)=\frac13x^3$ vì $F'(x)=x^2=f(x)$. Nhưng $G(x)=\frac13x^3+100$ cũng thỏa $G'(x)=x^2$. Vậy cả $F,G$ đều là nguyên hàm của $f$. Thật ra, mọi hàm $H(x)=\frac13x^3+C$, với $C$ là hằng số, đều là nguyên hàm. Câu hỏi đặt ra: còn những nguyên hàm nào khác không?

Để trả lời, nhắc lại ở §4.2 ta dùng định lý giá trị trung bình chứng minh hai hàm có cùng đạo hàm trên một khoảng chỉ sai khác nhau một hằng số (Hệ quả 4.2.7). Nếu $F,G$ là hai nguyên hàm bất kỳ của $f$, thì
\[F'(x)=f(x)=G'(x),\]
nên $G(x)-F(x)=C$, với $C$ là hằng số. Viết lại $G(x)=F(x)+C$, ta có kết quả sau.
\Needspace{7\baselineskip}
\group{1. Định lý}
Nếu $F$ là một nguyên hàm của $f$ trên khoảng $I$, thì nguyên hàm tổng quát nhất của $f$ trên $I$ là
\[F(x)+C,\]
trong đó $C$ là hằng số tùy ý.

Trở lại $f(x)=x^2$, nguyên hàm tổng quát là $\frac13x^3+C$. Gán các giá trị cụ thể cho $C$, ta được một họ hàm có các đồ thị là các phép tịnh tiến theo phương thẳng đứng của nhau (Hình 1). Điều này hợp lý vì các đường cong phải có cùng hệ số góc tại mỗi giá trị $x$ cho trước.
\fig[.46]{fig01}{Hình 1. Một số thành viên của họ nguyên hàm của $f(x)=x^2$.}

\Needspace{9\baselineskip}
\ex{1} Tìm nguyên hàm tổng quát nhất của mỗi hàm sau:

(a) $f(x)=\sin x$;\quad (b) $f(x)=1/x$;\quad (c) $f(x)=x^n$, $n\ne-1$.
\sol (a) Nếu $F(x)=-\cos x$ thì $F'(x)=\sin x$, nên $-\cos x$ là một nguyên hàm của $\sin x$. Theo Định lý 1, nguyên hàm tổng quát nhất là $G(x)=-\cos x+C$.

(b) Nhắc lại từ §3.6:
\[\frac d{dx}(\ln x)=\frac1x.\]
Vì vậy, trên $(0,\infty)$, nguyên hàm tổng quát của $1/x$ là $\ln x+C$. Ta cũng đã biết
\[\frac d{dx}(\ln|x|)=\frac1x\qquad(x\ne0).\]
Định lý 1 cho biết nguyên hàm tổng quát của $f(x)=1/x$ là $\ln|x|+C$ trên bất kỳ khoảng nào không chứa $0$. Đặc biệt, điều này đúng trên từng khoảng $(-\infty,0)$ và $(0,\infty)$. Do đó có thể viết nguyên hàm tổng quát của $f$ là
\[F(x)=\begin{cases}\ln x+C_1,&x>0,\\\ln(-x)+C_2,&x<0.\end{cases}\]
(c) Dùng quy tắc lũy thừa tìm một nguyên hàm của $x^n$. Nếu $n\ne-1$ thì
\[\frac d{dx}\left(\frac{x^{n+1}}{n+1}\right)=\frac{(n+1)x^n}{n+1}=x^n.\]
Vì vậy nguyên hàm tổng quát là
\[F(x)=\frac{x^{n+1}}{n+1}+C.\]
Công thức đúng với $n\ge0$ vì khi đó $f(x)=x^n$ xác định trên một khoảng. Nếu $n<0$, $n\ne-1$, công thức đúng trên mọi khoảng không chứa $0$ mà hàm được xác định.

\subsection*{Các công thức tìm nguyên hàm}
Như Ví dụ 1, mỗi công thức đạo hàm đọc từ phải sang trái cho một công thức nguyên hàm. Bảng 2 liệt kê một số nguyên hàm riêng. Mỗi công thức đúng vì đạo hàm của hàm ở cột phải là hàm ở cột trái. Công thức đầu cho biết nguyên hàm của hằng số nhân với một hàm bằng hằng số ấy nhân với nguyên hàm của hàm. Công thức thứ hai cho biết nguyên hàm của tổng bằng tổng các nguyên hàm. Ta dùng ký hiệu $F'=f$, $G'=g$.
\Needspace{19\baselineskip}
\group{2. Bảng công thức nguyên hàm}
\begin{center}\renewcommand{\arraystretch}{2.0}
\begin{tabular}{llll}\toprule
Hàm số&Nguyên hàm riêng&Hàm số&Nguyên hàm riêng\\\midrule
$cf(x)$&$cF(x)$&$\sin x$&$-\cos x$\\
$f(x)+g(x)$&$F(x)+G(x)$&$\sec^2x$&$\tan x$\\
$x^n\ (n\ne-1)$&$\dfrac{x^{n+1}}{n+1}$&$\sec x\tan x$&$\sec x$\\
$\dfrac1x$&$\ln|x|$&$\dfrac1{\sqrt{1-x^2}}$&$\arcsin x$\\
$e^x$&$e^x$&$\dfrac1{1+x^2}$&$\arctan x$\\
$b^x$&$\dfrac{b^x}{\ln b}$&$\cosh x$&$\sinh x$\\
$\cos x$&$\sin x$&$\sinh x$&$\cosh x$\\\bottomrule
\end{tabular}\end{center}
Để thu được nguyên hàm tổng quát nhất từ các nguyên hàm riêng trong Bảng 2, phải thêm một hằng số, hoặc các hằng số, như Ví dụ 1.

\Needspace{9\baselineskip}
\ex{2} Tìm mọi hàm $g$ thỏa
\[g'(x)=4\sin x+\frac{2x^5-\sqrt x}{x}.\]
\sol Trước hết viết lại
\begin{align*}
g'(x)&=4\sin x+\frac{2x^5}{x}-\frac{\sqrt x}{x}\\
&=4\sin x+2x^4-\frac1{\sqrt x}.
\end{align*}
Vậy cần tìm nguyên hàm của $g'(x)=4\sin x+2x^4-x^{-1/2}$. Dùng Bảng 2 và Định lý 1:
\begin{align*}
g(x)&=4(-\cos x)+2\frac{x^5}{5}-\frac{x^{1/2}}{1/2}+C\\
&=-4\cos x+\frac25x^5-2\sqrt x+C.
\end{align*}
Ta thường dùng chữ hoa $F$ để chỉ nguyên hàm của $f$. Nếu bắt đầu từ ký hiệu đạo hàm $f'$, thì tất nhiên một nguyên hàm là $f$.

Trong ứng dụng, thường gặp tình huống như Ví dụ 2: cần tìm hàm khi biết thông tin về các đạo hàm. Phương trình chứa đạo hàm của một hàm gọi là \textbf{phương trình vi phân}. Chúng sẽ được nghiên cứu kỹ hơn ở Chương 9; trước mắt, ta giải được một số phương trình vi phân sơ cấp. \textbf{Nghiệm tổng quát} chứa một hoặc nhiều hằng số tùy ý như Ví dụ 2. Tuy nhiên, các điều kiện bổ sung có thể xác định các hằng số, từ đó xác định duy nhất nghiệm.

\Needspace{9\baselineskip}
\ex{3} Tìm $f$ nếu $f'(x)=e^x+20(1+x^2)^{-1}$ và $f(0)=-2$.
\sol Nguyên hàm tổng quát của
\[f'(x)=e^x+\frac{20}{1+x^2}\]
là
\[f(x)=e^x+20\arctan x+C.\]
Để tìm $C$, dùng $f(0)=-2$:
\[f(0)=e^0+20\arctan0+C=-2.\]
Suy ra $C=-2-1=-3$, nên nghiệm riêng là
\[f(x)=e^x+20\arctan x-3.\]
\fig[.46]{fig02}{Hình 2}
Hình 2 cho đồ thị $f'$ và nguyên hàm $f$ trong Ví dụ 3. Vì $f'(x)>0$, $f$ luôn đồng biến. Khi $f'$ có cực đại hoặc cực tiểu, $f$ dường như có điểm uốn. Như vậy, đồ thị giúp kiểm tra phép tính.

\Needspace{9\baselineskip}
\ex{4} Tìm $f$ nếu $f''(x)=12x^2+6x-4$, $f(0)=4$, $f(1)=1$.
\sol Nguyên hàm tổng quát của $f''$ là
\[f'(x)=12\frac{x^3}{3}+6\frac{x^2}{2}-4x+C=4x^3+3x^2-4x+C.\]
Lấy nguyên hàm một lần nữa:
\begin{align*}
f(x)&=4\frac{x^4}{4}+3\frac{x^3}{3}-4\frac{x^2}{2}+Cx+D\\
&=x^4+x^3-2x^2+Cx+D.
\end{align*}
Dùng hai điều kiện để xác định $C,D$. Từ $f(0)=0+D=4$ có $D=4$. Vì
\[f(1)=1+1-2+C+4=1,\]
ta được $C=-3$. Vậy
\[f(x)=x^4+x^3-2x^2-3x+4.\]

\subsection*{Vẽ đồ thị nguyên hàm}
Nếu biết đồ thị $f$, có vẻ hợp lý rằng ta có thể phác họa đồ thị một nguyên hàm $F$. Chẳng hạn, biết $F(0)=1$, ta có điểm xuất phát $(0,1)$, và hướng di chuyển bút tại mỗi bước được cho bởi $F'(x)=f(x)$. Ví dụ tiếp dùng các nguyên tắc của chương này để vẽ $F$ ngay cả khi chưa biết công thức $f$. Điều đó xảy ra, chẳng hạn, khi $f(x)$ được xác định từ dữ liệu thực nghiệm.
\Needspace{8\baselineskip}
\ex{5} Đồ thị $f$ được cho ở Hình 3. Phác họa một nguyên hàm $F$ biết $F(0)=2$.
\fig[.46]{fig03}{Hình 3}
\sol Dùng tính chất hệ số góc của $y=F(x)$ bằng $f(x)$. Bắt đầu tại $(0,2)$, ta vẽ $F$ ban đầu nghịch biến vì $f(x)<0$ khi $0<x<1$. Do $f(1)=f(3)=0$, $F$ có tiếp tuyến ngang tại $x=1,3$. Với $1<x<3$, $f(x)>0$, nên $F$ đồng biến. $F$ có cực tiểu địa phương tại $x=1$ và cực đại địa phương tại $x=3$. Với $x>3$, $f(x)<0$, nên $F$ nghịch biến trên $(3,\infty)$. Vì $f(x)\to0$ khi $x\to\infty$, đồ thị $F$ ngày càng phẳng khi $x\to\infty$. Ngoài ra, $F''(x)=f'(x)$ đổi dấu từ dương sang âm tại $x=2$, từ âm sang dương tại $x=4$, nên $F$ có điểm uốn tại $x=2,4$. Dùng các thông tin này, ta phác họa đồ thị ở Hình 4.
\fig[.46]{fig04}{Hình 4}

\subsection*{Chuyển động thẳng}
Phép tìm nguyên hàm đặc biệt hữu ích khi phân tích chuyển động trên đường thẳng. Nếu hàm vị trí là $s=f(t)$, hàm vận tốc là $v(t)=s'(t)$. Vậy hàm vị trí là một nguyên hàm của vận tốc. Tương tự, $a(t)=v'(t)$, nên vận tốc là một nguyên hàm của gia tốc. Nếu biết gia tốc và các giá trị ban đầu $s(0),v(0)$, có thể tìm vị trí bằng hai lần lấy nguyên hàm.

\Needspace{9\baselineskip}
\ex{6} Chất điểm chuyển động thẳng với gia tốc $a(t)=6t+4$, vận tốc ban đầu $v(0)=-6$ cm/s và độ dời ban đầu $s(0)=9$ cm. Tìm hàm vị trí $s(t)$.
\sol Vì $v'(t)=a(t)=6t+4$, lấy nguyên hàm:
\[v(t)=6\frac{t^2}{2}+4t+C=3t^2+4t+C.\]
Ta có $v(0)=C=-6$, nên
\[v(t)=3t^2+4t-6.\]
Vì $s'=v$, lấy nguyên hàm một lần nữa:
\[s(t)=3\frac{t^3}{3}+4\frac{t^2}{2}-6t+D=t^3+2t^2-6t+D.\]
Điều này cho $s(0)=D=9$, nên hàm vị trí cần tìm là
\[s(t)=t^3+2t^2-6t+9.\]

Vật gần mặt đất chịu trọng lực gây gia tốc hướng xuống, ký hiệu $g$. Với chuyển động gần mặt đất, có thể coi $g$ không đổi, khoảng $9{,}8$ m/s$^2$ hoặc $32$ ft/s$^2$. Đáng chú ý là chỉ từ tính không đổi của gia tốc trọng trường, ta dùng phép tính vi phân để suy ra vị trí và vận tốc của vật chuyển động dưới tác dụng trọng lực, như ví dụ tiếp theo.

\Needspace{9\baselineskip}
\par\medskip\noindent\begin{minipage}{\linewidth}
\ex{7} Quả bóng được ném lên với tốc độ 48 ft/s từ mép vách đá cao 432 ft so với mặt đất. Tìm độ cao sau $t$ giây. Khi nào bóng đạt độ cao lớn nhất? Khi nào bóng chạm đất?
\sol Chuyển động theo phương thẳng đứng, chọn chiều dương hướng lên. Tại $t$, độ cao là $s(t)$ và vận tốc $v(t)$ giảm. Gia tốc âm:
\[a(t)=\frac{dv}{dt}=-32.\]
Lấy nguyên hàm, $v(t)=-32t+C$. Dùng $v(0)=48$ có $48=0+C$, nên
\[v(t)=-32t+48.\]
Bóng đạt độ cao lớn nhất khi $v(t)=0$, tức sau 1,5 giây. Vì $s'=v$, lấy nguyên hàm lần nữa:
\[s(t)=-16t^2+48t+D.\]
Dùng $s(0)=432$ có $D=432$, nên
\[s(t)=-16t^2+48t+432.\]
Biểu thức đúng đến khi bóng chạm đất. Điều đó xảy ra khi $s(t)=0$:
\[-16t^2+48t+432=0,\]
tương đương
\[t^2-3t-27=0.\]
Công thức nghiệm phương trình bậc hai cho
\[t=\frac{3\pm3\sqrt{13}}2.\]
Loại nghiệm mang dấu trừ vì $t<0$. Vậy bóng chạm đất sau $3(1+\sqrt{13})/2\approx6{,}9$ giây.
\fig[.46]{fig05}{Hình 5}
Đồ thị hàm vị trí ở Hình 5 xác nhận các kết luận: bóng đạt độ cao lớn nhất sau 1,5 giây và chạm đất sau khoảng 6,9 giây.

\end{minipage}\par
\clearpage
\section*{4.9 — Bài tập}
\Needspace{9\baselineskip}
\group{1–4}
Tìm một nguyên hàm của hàm số.
\bt{1} (a) $f(x)=6$;\quad (b) $g(t)=3t^2$.
\bt{2} (a) $f(x)=2x$;\quad (b) $g(x)=-1/x^2$.
\bt{3} (a) $h(q)=\cos q$;\quad (b) $f(x)=e^x$.
\bt{4} (a) $g(t)=1/t$;\quad (b) $r(\theta)=\sec^2\theta$.
\Needspace{10\baselineskip}
\group{5–26}
Tìm nguyên hàm tổng quát nhất. Kiểm tra bằng cách lấy đạo hàm kết quả.
\bt{5} $f(x)=4x+7$.
\bt{6} $f(x)=x^2-3x+2$.
\bt{7} $f(x)=2x^3-\frac23x^2+5x$.
\bt{8} $f(x)=6x^5-8x^4-9x^2$.
\bt{9} $f(x)=x(12x+8)$.
\bt{10} $f(x)=(x-5)^2$.
\bt{11} $g(x)=4x^{-2/3}-2x^{5/3}$.
\bt{12} $h(z)=3z^{0{,}8}+z^{-2{,}5}$.
\bt{13} $f(x)=3\sqrt x-2\sqrt[3]x$.
\bt{14} $g(x)=\sqrt x(2-x+6x^2)$.
\bt{15} $\displaystyle f(t)=\frac{2t-4+3\sqrt t}{\sqrt t}$.
\bt{16} $f(x)=\sqrt[4]5+\sqrt[4]x$.
\bt{17} $\displaystyle f(x)=\frac2{5x}-\frac3{x^2}$.
\bt{18} $\displaystyle f(x)=\frac{5x^2-6x+4}{x^2}$, $x>0$.
\bt{19} $g(t)=7e^t-e^3$.
\bt{20} $\displaystyle f(x)=\frac{10}{x^6}-2e^x+3$.
\bt{21} $f(\theta)=2\sin\theta-3\sec\theta\tan\theta$.
\bt{22} $h(x)=\sec^2x+\pi\cos x$.
\bt{23} $\displaystyle f(r)=\frac4{1+r^2}-\sqrt[5]{r^4}$.
\bt{24} $\displaystyle g(v)=2\cos v-\frac3{\sqrt{1-v^2}}$.
\bt{25} $f(x)=2^x+4\sinh x$.
\bt{26} $\displaystyle f(x)=\frac{2x^2+5}{x^2+1}$.
\Needspace{10\baselineskip}
\group{27–28}
Tìm nguyên hàm $F$ của $f$ thỏa điều kiện đã cho. Kiểm tra kết quả bằng cách so sánh đồ thị $f$ và $F$.
\bt{27} $f(x)=2e^x-6x$,\quad $F(0)=1$.
\bt{28} $f(x)=4-3(1+x^2)^{-1}$,\quad $F(1)=0$.
\Needspace{8\baselineskip}
\group{29–54}
Tìm $f$.
\bt{29} $f''(x)=24x$.
\bt{30} $f''(t)=t^2-4$.
\bt{31} $f''(x)=4x^3+24x-1$.
\bt{32} $f''(x)=6x-x^4+3x^5$.
\bt{33} $f''(x)=2x+3e^x$.
\bt{34} $f''(x)=1/x^2$,\quad $x>0$.
\bt{35} $f'''(t)=12+\sin t$.
\bt{36} $f'''(t)=\sqrt t-2\cos t$.
\bt{37} $f'(x)=8x^3+1/x$,\quad $x>0$,\quad $f(1)=-3$.
\bt{38} $f'(x)=\sqrt x-2$,\quad $f(9)=4$.
\bt{39} $f'(t)=4/(1+t^2)$,\quad $f(1)=0$.
\bt{40} $f'(t)=t+1/t^3$,\quad $t>0$,\quad $f(1)=6$.
\bt{41} $f'(x)=5x^{2/3}$,\quad $f(8)=21$.
\bt{42} $f'(x)=(x+1)/\sqrt x$,\quad $f(1)=5$.
\bt{43} $f'(t)=\sec t(\sec t+\tan t)$,\quad $-\pi/2<t<\pi/2$,\quad $f(\pi/4)=-1$.
\bt{44} $f'(t)=3t^2-3/t$,\quad $f(1)=2$,\quad $f(-1)=1$.
\bt{45} $f''(x)=-2+12x-12x^2$,\quad $f(0)=4$,\quad $f'(0)=12$.
\bt{46} $f''(x)=8x^3+5$,\quad $f(1)=0$,\quad $f'(1)=8$.
\bt{47} $f''(\theta)=\sin\theta+\cos\theta$,\quad $f(0)=3$,\quad $f'(0)=4$.
\bt{48} $f''(t)=t^2+1/t^2$,\quad $t>0$,\quad $f(2)=3$,\quad $f'(1)=2$.
\bt{49} $f''(x)=4+6x+24x^2$,\quad $f(0)=3$,\quad $f(1)=10$.
\bt{50} $f''(x)=x^3+\sinh x$,\quad $f(0)=1$,\quad $f(2)=2{,}6$.
\bt{51} $f''(x)=e^x-2\sin x$,\quad $f(0)=3$,\quad $f(\pi/2)=0$.
\bt{52} $f''(t)=\sqrt[3]t-\cos t$,\quad $f(0)=2$,\quad $f(1)=2$.
\bt{53} $f''(x)=x^{-2}$,\quad $x>0$,\quad $f(1)=0$,\quad $f(2)=0$.
\bt{54} $f'''(x)=\cos x$,\quad $f(0)=1$,\quad $f'(0)=2$,\quad $f''(0)=3$.
\bt{55} Đồ thị $f$ đi qua $(2,5)$ và hệ số góc tiếp tuyến tại $(x,f(x))$ là $3-4x$. Tìm $f(1)$.
\bt{56} Tìm hàm $f$ sao cho $f'(x)=x^3$ và đường thẳng $x+y=0$ tiếp xúc đồ thị $f$.
\Needspace{11\baselineskip}
\group{57–58}
Đồ thị $f$ được cho trong hình. Đồ thị nào là một nguyên hàm của $f$? Vì sao?
\par\medskip\noindent\begin{minipage}{\linewidth}
\bt{57} Xét hình sau.
\fig[.46]{ex57}{Hình cho Bài tập 57.}
\end{minipage}\par
\par\medskip\noindent\begin{minipage}{\linewidth}
\bt{58} Xét hình sau.
\fig[.46]{ex58}{Hình cho Bài tập 58.}
\end{minipage}\par
\par\medskip\noindent\begin{minipage}{\linewidth}
\bt{59} Đồ thị hàm số được cho trong hình. Phác họa một nguyên hàm $F$, biết $F(0)=1$.
\fig[.50]{ex59}{Hình cho Bài tập 59.}
\end{minipage}\par
\par\medskip\noindent\begin{minipage}{\linewidth}
\bt{60} Hình vẽ cho đồ thị vận tốc của một chất điểm. Phác họa đồ thị một hàm vị trí.
\fig[.50]{ex60}{Hình cho Bài tập 60.}
\end{minipage}\par
\par\medskip\noindent\begin{minipage}{\linewidth}
\bt{61} Hình vẽ cho đồ thị $f'$. Phác họa đồ thị $f$ nếu $f$ liên tục trên $[0,3]$ và $f(0)=-1$.
\fig[.43]{ex61}{Hình cho Bài tập 61.}
\end{minipage}\par
\bt{62} (a) Vẽ $f(x)=2x-3\sqrt x$.

(b) Từ đồ thị ở (a), phác họa nguyên hàm $F$ thỏa $F(0)=1$.

(c) Dùng các quy tắc trong mục này tìm biểu thức $F(x)$.

(d) Vẽ $F$ từ biểu thức ở (c), rồi so sánh với hình phác họa ở (b).
\Needspace{9\baselineskip}
\group{63–64}
Vẽ đồ thị $f$ và dùng nó để phác họa nguyên hàm đi qua gốc tọa độ.
\bt{63} $\displaystyle f(x)=\frac{\sin x}{1+x^2}$,\quad $-2\pi\le x\le2\pi$.
\bt{64} $f(x)=\sqrt{x^4-2x^2+2}-2$,\quad $-3\le x\le3$.
\Needspace{9\baselineskip}
\group{65–70}
Chất điểm chuyển động với các dữ kiện sau. Tìm vị trí chất điểm.
\bt{65} $v(t)=2\cos t+4\sin t$,\quad $s(0)=3$.
\bt{66} $v(t)=t^2-3\sqrt t$,\quad $s(4)=8$.
\bt{67} $a(t)=2t+1$,\quad $s(0)=3$,\quad $v(0)=-2$.
\bt{68} $a(t)=3\cos t-2\sin t$,\quad $s(0)=0$,\quad $v(0)=4$.
\bt{69} $a(t)=\sin t-\cos t$,\quad $s(0)=0$,\quad $s(\pi)=6$.
\bt{70} $a(t)=t^2-4t+6$,\quad $s(0)=0$,\quad $s(1)=20$.
\bt{71} Thả một hòn đá từ đài quan sát trên cao (Space Deck) của tháp CN, cao 450 m so với mặt đất.

(a) Tìm độ cao của đá so với mặt đất tại thời điểm $t$.

(b) Đá mất bao lâu để chạm đất?

(c) Vận tốc khi chạm đất bằng bao nhiêu?

(d) Nếu ném đá xuống với tốc độ 5 m/s, mất bao lâu để chạm đất?
\bt{72} Chứng minh với chuyển động thẳng có gia tốc không đổi $a$, vận tốc ban đầu $v_0$, độ dời ban đầu $s_0$, độ dời sau thời gian $t$ là $s=\frac12at^2+v_0t+s_0$.
\bt{73} Một vật được phóng lên với vận tốc ban đầu $v_0$ m/s từ vị trí cao $s_0$ mét so với mặt đất. Chứng minh
\[[v(t)]^2=v_0^2-19{,}6[s(t)-s_0].\]
\bt{74} Hai quả bóng được ném lên từ mép vách đá ở Ví dụ 7. Quả thứ nhất có tốc độ 48 ft/s; quả thứ hai được ném sau một giây với tốc độ 24 ft/s. Hai quả bóng có bao giờ đi ngang qua nhau không?
\bt{75} Hòn đá được thả từ vách đá, chạm đất với tốc độ 120 ft/s. Vách đá cao bao nhiêu?
\par\medskip\noindent\begin{minipage}{\linewidth}
\bt{76} Người nhảy cầu có khối lượng $m$ đứng ở đầu ván dài $L$, khối lượng trên đơn vị chiều dài $\rho$. Ván có dạng $y=f(x)$, với
\[EIy''=mg(L-x)+\frac12\rho g(L-x)^2,\]
trong đó $E,I>0$ là các hằng số phụ thuộc vật liệu ván, còn $g<0$ là gia tốc trọng trường.

(a) Tìm biểu thức mô tả hình dạng đường cong.

(b) Dùng $f(L)$ ước lượng độ võng xuống dưới phương ngang ở đầu ván.
\fig[.50]{ex76}{Hình cho Bài tập 76.}
\end{minipage}\par
\bt{77} Công ty ước lượng chi phí biên, tính bằng USD/sản phẩm, khi sản xuất $x$ sản phẩm là $1{,}92-0{,}002x$. Chi phí sản xuất một sản phẩm là 562 USD. Tìm chi phí sản xuất 100 sản phẩm.
\bt{78} Thanh dài 1 m có khối lượng trên đơn vị chiều dài $\rho(x)=1/\sqrt x$ g/cm, trong đó $x$ tính bằng cm từ một đầu thanh. Tìm khối lượng thanh.
\Needspace{13\baselineskip}
\bt{79} Giọt mưa lớn dần khi rơi, diện tích bề mặt và lực cản vì thế tăng. Vận tốc ban đầu hướng xuống là 10 m/s, gia tốc hướng xuống là
\[a=\begin{cases}9-0{,}9t,&0\le t\le10,\\0,&t>10.\end{cases}\]
Nếu giọt mưa hình thành ở độ cao 500 m, nó mất bao lâu để rơi xuống đất?
\bt{80} Ô tô đang chạy 50 mi/h thì phanh hết mức, tạo độ giảm tốc không đổi 22 ft/s$^2$. Xe đi thêm bao xa trước khi dừng?
\bt{81} Cần gia tốc không đổi bằng bao nhiêu để tăng tốc ô tô từ 30 mi/h lên 50 mi/h trong 5 giây?
\bt{82} Ô tô phanh với độ giảm tốc không đổi 16 ft/s$^2$, để lại vết trượt dài 200 ft trước khi dừng. Tốc độ lúc bắt đầu phanh là bao nhiêu?
\bt{83} Xe chạy 100 km/h, người lái thấy tai nạn cách 80 m phía trước và đạp phanh gấp. Cần độ giảm tốc không đổi bao nhiêu để dừng kịp, tránh va chạm liên hoàn?
\bt{84} Tên lửa mô hình được phóng thẳng đứng từ trạng thái nghỉ. Trong ba giây đầu, gia tốc $a(t)=60t$; sau đó hết nhiên liệu và tên lửa trở thành vật rơi tự do. Mười bốn giây sau, dù mở và tốc độ rơi giảm tuyến tính trong 5 giây, đến khi vận tốc bằng $-18$ ft/s. Sau đó tên lửa tiếp tục hạ xuống với vận tốc này.

(a) Xác định hàm vị trí $s$, vận tốc $v$ cho mọi thời điểm $t$. Phác họa đồ thị $s,v$.

(b) Tên lửa đạt độ cao lớn nhất lúc nào, và độ cao ấy là bao nhiêu?

(c) Tên lửa chạm đất lúc nào?
\bt{85} Một tàu cao tốc tăng tốc và giảm tốc với độ lớn 2,4 ft/s$^2$. Tốc độ hành trình lớn nhất là 180 mi/h.

(a) Quãng đường lớn nhất tàu đi được nếu tăng tốc từ nghỉ đến tốc độ hành trình rồi chạy ở tốc độ đó trong 20 phút là bao nhiêu?

(b) Giả sử tàu xuất phát từ nghỉ và phải dừng hoàn toàn sau 20 phút. Quãng đường lớn nhất đi được là bao nhiêu?

(c) Tìm thời gian ngắn nhất tàu đi giữa hai ga liên tiếp cách nhau 60 dặm.

(d) Chuyến đi giữa hai ga mất 37,5 phút. Hai ga cách nhau bao xa?
\end{document}
